\songtitle{Lost in What}

% -> personal.tex (Biographical); stories.tex (Self); pop.tex (K-pop & J-pop)
\tag*{2026}\tag{j-pop}\tag{original-source}
\begin{notes}
Lyrics by SUNO based on a short story extracted from Carlos Thompson's memoir \emph{A Life in World Events}, cowritten with Claude -- a bronze medal at a Beijing math olympiad, a solo climb up Mount Fuji, and the jet-lagged disorientation of a week spanning sixteen time zones between Bogotá and his parents' home in Stockholm. Retrieved from Wyomee Dank's Suno page (V6-MINI, 26 September 2026).
\end{notes}
\begin{style}
J-pop with bright, finely articulated piano leading tight modern pop drums, rounded bass, and clean electric guitar that widen the choruses; two distinct shakuhachi-and-koto instrumental bridges over restrained taiko, then a spacious pause before the final chorus. Polished contemporary production, vivid controlled dynamics, measured mid-tempo groove; clear earnest male tenor, close in the verses with layered final-chorus harmonies.
\end{style}
\begin{lyrics}
\songpart{Verse 1}
Beijing heat, the papers turned\\
Third time beneath that contest clock\\
Bronze again, two points away\\
From silver's line in Germany\\
Tokyo gave us borrowed nights\\
Then Fuji called before the dawn\\
My stomach failed; my friends moved on\\
Their footsteps vanished up the slope

\songpart{Pre-Chorus}
I waited for the sky to pale\\
Then took the path alone

\songpart{Chorus}
I knew the mountain, knew my place\\
Could trace the summit, name the trail\\
But every station down below\\
Was one more answer I couldn't know\\
Not lost in where, but lost in what\\
A perfect map could not unlock\\
I knew exactly where I was\\
Still had no way to choose the road

\songpart{Verse 2}
The morning opened, cold and wide\\
I reached the top with no one near\\
Then turned to make the long descent\\
With every bus stop drawing near\\
I held the thought against the climb:\\
A point can fix you on the ground\\
And still leave every next step blank\\
While all the facts are true around

\annotation{Traditional Japanese Bridge}
\annotation{Instrumental: shakuhachi and koto with restrained taiko}

\songpart{Chorus}
I knew the mountain, knew my place\\
Could trace the summit, name the trail\\
But every station down below\\
Was one more answer I couldn't know\\
Not lost in where, but lost in what\\
A perfect map could not unlock\\
I knew exactly where I was\\
Still had no way to choose the road

\songpart{Verse 3}
At last I found my team again\\
Together back to Tokyo\\
Bogotá nights, then north alone\\
To Stockholm, where my parents lived\\
Three in the morning, wide awake\\
The window held a noon-bright sky\\
Sixteen time zones in a week\\
One day slipped past the date line

\songpart{Pre-Chorus}
I knew the room, the city, home\\
My body hadn't caught up yet

\annotation{Traditional Japanese Bridge}
\annotation{Instrumental: koto and shakuhachi, rising into a full-band transition}

\songpart{Final Chorus}
I knew the window, knew my place\\
Could name the city, read the sky\\
But daylight found me deep in night\\
And all my bearings disagreed inside\\
Not lost in where, but lost in what\\
A perfect map could not unlock\\
I knew exactly where I was\\
Still had no way to choose the road

\songpart{Outro}
We found our way back, side by side\\
My body took the longer ride
\end{lyrics}

